\documentclass[11pt, oneside]{book}
\usepackage[spanish]{babel}
\usepackage[margin=1in]{geometry}
\usepackage{tikz}
\usepackage{soul}
\usepackage{float}
\usepackage{wrapfig}
\usepackage{amsmath}
\usepackage{multicol}
\usepackage{sectsty}
\usepackage[utf8]{inputenc}
% \usepackage{mhchem}
\usepackage{enumitem}
\usetikzlibrary{decorations.markings} 

\definecolor{myyellow}{RGB}{254,241,24}
\definecolor{myorange}{RGB}{234,125,1}

\restylefloat{figure}



\allsectionsfont{\flushleft}
\begin{document}

\title{CURSO DE PROTECCI\'ON RADIOL\'OGICA PARA OPERADORES DE FUENTES PARA
  RADIOGRAF\'IA INDUSTRIAL}
\maketitle

\tikz [remember picture, overlay] %
\node [shift={(1cm,-1cm)}] at (current page.north west) %
[anchor=north west] %
{
\includegraphics{logo.svg}};
\thispagestyle{empty}

\chapter*{OBJETIVOS Y PROGRAMA}

\section*{Objetivos}
Adiestrar al operador para el manejo seguro de las fuentes de
radiación empleadas en radiografía.

El Operador:
\begin{enumerate}
  \item Conocerá los efectos biológicos de la radiación y los principios
    para la limitación de dosis.
  \item Aplicara los principios básicos de protección contra las
    radiaciones.
  \item Manejara los instrumentos para medición de radiaciones, así como
    los dosímetros personales y alarmas sonoras.
  \item Incrementara su sentido de responsabilidad con relación a su
    seguridad personal, la de sus compañeros y la del publico en
    general.
  \item Manejara adecuadamente los incidentes que frecuentemente se
    presenten en el trabajo y que es factible resolver con los medios a
    su alcance.
  \item Distinguirá cuando debe resolver una situación y cuando debe
    pedir ayuda para su solución.
\end{enumerate}

\section*{Programa}
\begin{enumerate}
  \item \ul{MOTIVACI\'ON} \hfill (1 hr)

    Motivar al grupo respecto a la importancia de la protección
    radiológica y de la necesidad de que seguridad constituya un habito
    personal. Explicara la responsabilidad adquirida al manejar una
    fuente radioactiva.
  \item \ul{ESTRUCTURA AT\'OMICA (TERMINOLOG\'IA)} \hfill (1 hr)

    El alumno identificara definiciones de: átomo, núcleo, protón,
    electrón, numero atómico, numero de masa, numero de neutrones,
    aplicara correctamente esas definiciones y la simbología ${}^A_ZX$.
  \item \ul{RADIOACTIVIDAD Y RADIACI\'ON} \hfill (1:30 hr)

    El alumno podrá describir el fenómeno de la desintegración beta y
    gamma. Explicara el concepto de vida media y podrá determinar la
    actividad actual de una fuente conociendo su vida media, actividad
    inicial y el tiempo transcurrido, empleando una gráfica o una tabla
    de decaimiento. Empleara correctamente las unidades de actividad:
    Curie Y Becquerel.

  \item \ul{INTERACCI\'ON DE LA RADIACI\'ON CON LA MATERIA} \hfill (1
    hr)

    Definirá con su palabras cada uno de los siguientes conceptos: ion,
    ionización, excitación, absorción, alcance o penetración,
    dispersión, radiación primaria, radiación dispersa y radiación
    trasmitida.

    \textbf{PRACTICA NO. 1} \ul{Absorci\'on de las radiaciones beta
    y gamma} \hfill (1 hr)

  \item \ul{EXPOSICI\'ON, DOSIS Y UNIDADES} \hfill (1 hr)

    El alumno explicara la diferencia entre exposición y dosis,
    exposición y rapidez de exposición. Definirá Roentgen, el rad y sus
    submúltiplos, así como las unidades de rapidez de exposición y de
    dosis. Describirá los conceptos de equivalente de dosis y rem, así
    como las unidades Gray y Sievert y sus equivalentes a rad y rem.

  \item \ul{EFECTOS BIOL\'OGICOS DE LA RADIACI\'ON} \hfill (2:30 hr)
    \begin{enumerate}[label*=\arabic*.]
      \item \ul{Efectos Som\'aticos y Efectos Gen\'eticos}

        Podrá señalar la diferencia entre los efectos somáticos y los
        efectos genéticos, así como la diferente radio-sensibilidad de los
        distintos tipos de células, tejidos y órganos.
      \item \ul{Relaci\'on Dosis-Efecto}

        Describirá la diferencia entre dosis aguda y dosis crónica,
        distinguirá los valores de dosis que producen los distintos
        efectos, desde un eritem hasta una dosis letal, así como la
        diferencia entre una irradiación localizada y una irradiación a
        cuerpo entero.
      \item \ul{Formas de Exposici\'on a la Radiaci\'on}

        El alumno distinguirá entre casos de exposición a radiación
        externa, ingestión, inhalación, y dosis interna.
    \end{enumerate}
  \item \ul{MEDICI\'ON DE LA RADIACI\'ON} \hfill (0:30 hr)

    El alumno manejara los diversos instrumentos de monitoreo, escogerá
    el adecuado para cada situación e interpretara correctamente las
    lecturas: quedara convencido de la ventaja de usar siempre un
    dosímetro personal, una alarma sonora y un instrumento de monitoreo,
    así como de cuidarlos adecuadamente.
    \begin{enumerate}[label*=\arabic*.]
      \item \ul{Dispositivos de Vigilancia Personal}
        \begin{enumerate}
          \item \ul{Dos\'imetros de Bolsillo} \hfill (0:30 hr)

            Entenderá el funcionamiento y manejara un dosímetro de bolsillo
            y su cargador, podrá verificar su funcionamiento y tomar las
            medidas pertinentes para su conservación. Interpretara
            correctamente las lecturas y conocerá sus ventajas y
            limitaciones.

            \textbf{PRACTICA NO. 2} \ul{Manejo y uso del dos\'imetro}
            \hfill (1 hr)
          \item \ul{Dos\'imetros de Pel\'icula y Termo-luminiscentes}
            \hfill (0:30 hr)

            El alumno reconocerá y podrá explicar la importancia del uso de
            un dosímetro que registre la exposición acumulada. Se percatara
            del a necesidad de adquirir el habito de usarlo y cuidar de que
            su uso y conservación sean adecuados.

          \item \ul{Detector con Alarma Sonora} \hfill (0:30 hr)

            El alumno usara un detector con alarma sonora, relacionara la
            señal audible con el índice de exposición correspondiente y
            reconocerá las ventajas de su uso continuo.

            \textbf{PRACTICA NO. 3} \ul{Manejo del detector con
            alarma sonora} \hfill (1 hr)
        \end{enumerate}
      \item \ul{Instrumentos de Monitoreo}

        El alumno manejara correctamente al menos uno de los siguientes
        instrumentos de monitoreo con cámara de ionización, con detector
        Geiger o con detector de centello. Interpretara correctamente las
        lecturas tomando en consideración las limitaciones de cada
        uno. Podrá verificar la calibración del instrumento usando una
        fuente de referencia.
        \begin{enumerate}
          \item \ul{Monitor con detector Geiger} \hfill (1 hr)

            Lo manejara y explicara en que casos es conveniente su
            uso. Describirá sus ventajas y desventajas frente a una cámara
            de ionización. Describirá el fenómeno de saturación y el uso del
            factor de calibración.

            \textbf{PRACTICA NO. 4} \ul{Manejo de un monitor con
            detector Geiger} \hfill (1 hr)

            \textbf{PRACTICA NO. 5} \ul{Localizaci\'on de fuentes con
            Geiger-Muller} \hfill (1:30 hr)
        \end{enumerate}
    \end{enumerate}
  \item \ul{FILOSOF\'IA DE LA PROTECCI\'ON RADIOL\'OGICA} \hfill (2 hr)
    \begin{enumerate}[label*=\arabic*.]
      \item \ul{Concepto de Riesgo-Beneficio. Concepto de ALARA}

        El alumno podrá percatarse de las diferencias entre el riesgo y el
        beneficio, en diversos tipos de actividades y, en especial, en la
        radiografía industrial. Explicara en forma simplificada el
        concepto ALARA.

      \item \ul{Sistema de Limitaci\'on de Dosis}

        Memorizara el valor del limite de equivalente de dosis para
        trabajadores ocupacionalmente expuestos y los valore de rapidez de
        exposición en mr/hr y mr/semana que en promedio producirán dicho
        limite.
    \end{enumerate}
  \item \ul{PRINCIPIOS B\'ASICOS DE PROTECCI\'ON CONTRA RADIACI\'ON
    EXTERNA} \hfill (3:30 hr)
    \begin{enumerate}[label*=\arabic*.]
      \item \ul{Distancia y Tiempo Como Medios de Protecci\'on}
        \hfill (2 hr)

        El alumno calculara la rapidez de exposición a diferentes
        distancias de una fuente de actividad conocida. Calculara la
        exposición en términos de rapidez de exposición y el
        tiempo. Calculara el tiempo que puede permanecer una persona en un
        campo de radiación conocido, para recibir una exposición o una
        dosis predeterminada a distintas distancias de la fuente. Deberá
        quedar suficientemente familiarizado con este tipo de cálculos
        como para hacerlos sin dificultad en forma rutinaria.

        \textbf{PRACTICA NO. 6} \ul{Aplicaci\'on de la ley del cuadrado inverso.} \hfill (1:30 hr)

      \item \ul{El Blindaje Como Medio de Protecci\'on} \hfill (1 hr)

        El alumno podrá determinar la atenuación de un haz de radiación gamma
        al atravesar un blindaje de espesor y material conocidos empleando
        para ello una gráfica semilogaritmica y aplicando el concepto de
        una capa hemirreductora y capa decirreductora para determinar espesores
        de blindaje suficientes, a fin de producir la atenuación deseada.

        \textbf{PRACTICA NO. 7} \ul{Atenuaci\'on de la radiaci\'on gamma,determinaci\'on de la capa hemirreductora utilizando diversos materiales.}
        \hfill\mbox{(1 hr)}

      \item \ul{Protecci\'on Contra la Contaminaci\'on Externa y la Ingesti\'on o la Inhalaci\'on} \hfill\mbox{(0:30 hr)}

        El alumno reafirmara sus conceptos de contaminación y podrá
        explicar la diferencia entre contaminación externa, ingestión
        e inhalación, así como la forma en que pueden producirse y los
        medios para evitarla.

      \item \ul{La Protecci\'on Radiol\'ogica Durante el Trabajo} \hfill\mbox{(1 hr)}

        Identificara las situaciones de mayor riesgo durante el trabajo radiográfico.
        Describirá secuencialmente las medidas de protección que debe tomarse en
        una jornada de trabajo.
    \end{enumerate}

  \item \ul{EQUIPOS PARA RADIOGRAF\'IA INDUSTRIAL QUE UTILIZAN MATERIAL RADIOACTIVO}
    \begin{enumerate}[label*=\arabic*.]
      \item \ul{Principios de la Utilizaci\'on de la Fuente y Descripci\'on de los Equipos Asociados con el Uso.}\hfill\mbox{(1 hr)}

      \item \ul{Procedimientos Generales para el Transporte y Almacenamiento de Material Radioactivo.}\hfill\mbox{(1 hr)}

        El alumno reconocerá y podrá explicar la importancia de las medidas de Seguridad Física
        y Radiológica que deben contemplarse en el transporte y durante el almacén del material radioactivo.
        Describirá las posibles consecuencias de pasar por alto dichas medidas.
    \end{enumerate}
  \item \ul{NORMAS Y REGLAMENTOS}\hfill\mbox{(1:30 hr)}

    \ul{Objetivo General}: Conocer las principales disposiciones reglamentarias aplicables y las funciones del Organismo Regulador en
    materia de seguridad radiológica.
    \begin{enumerate}[label*=\arabic*.]
      \item \ul{Manual de Seguridad Radiol\'ogica.}

        El alumno dará una interpretación correcta (que le permita en el futuro aplicar adecuadamente
        a cada uno de las disposiciones del Manual Interno de Seguridad Radiológica. Distinguirá el contenido mínimo de un Manual
        de seguridad Radiológica para Radiografía Industrial.
      \item \ul{Estructura, Organizaci\'on y Cadena de Responsabilidad en Materia de Seguridad Radiol\'ogica.}

        El alumno podrá canalizar adecuadamente sus necesidades de equipo de seguridad,
        información y ayuda en caso de incidente, dentro de la estructura jerárquica de la empresa.

      \item \ul{Obligaciones del Encargado de Seguridad Radiol\'ogica.}

        El alumno describirá las funciones mas importantes del Encargado de Seguridad Radiológica.

      \item \ul{La C.N.S.N.S. sus Funciones y Organizaci\'on.}

        El alumno explicara cuales son las funciones y las dependencias
        de la C.N.S.N.S. relacionadas con su área de trabajo, así como la
        forma de acceso a esas dependencias cuando requiera de orientación
        o auxilio en caso de incidentes o accidentes radiológicos.

    \end{enumerate}

  \item \ul{AN\'ALISIS DE INCIDENTES Y ACCIDENTES (casos reales)}\hfill\mbox{(3 hr)}

    \ul{Objetivo General}: Motivar mediante la exposición de incidentes
    y accidentes reales, el análisis de las situaciones para buscar con
    calma y razonamiento lógico la solución que de lugar a la exposición
    mínima, evitando así acciones precipitadas que pueden conducir a riesgos
    mayores.

    \begin{enumerate}[label*=\arabic*.]
      \item \ul{Procedimiento de Emergencia}

        El alumno explicara como proceder en cada uno de
        los tipos de incidentes mas comunes, en los que
        este a su alcance la posibilidad de solución.
        Indicara a quien acudir cuando no disponga de los medios para
        resolver el caso, señalando las medidas o disposiciones
        inmediatas que debe tomar para evitar que se agrave la situación.

      \item \ul{Informe de Incidentes y Accidentes}

        El alumno indicara los tipos de de incidentes que deberá informar a sus superiores y la forma en que lo hará.

      \item \ul{Plan de Emergencia}

    \end{enumerate}
    \textbf{PRACTICA NO. 8} \ul{Manejo de un equipo radiogr\'afico, soluci\'on de fallas}\hfill\mbox{(2 hr)}

    \textbf{PRACTICA NO. 9} \ul{Recuperaci\'on de fuentes, uso de manipuladores a distancia}\hfill\mbox{(2 hr)}

    \textbf{PRACTICA NO. 9} \ul{Simulacro de un accidente}\hfill\mbox{(4 hr)}
\end{enumerate}

\begin{flushleft}
  EVALUACI\'ON FINAL TE\'ORICO-PRACTICO.

  Duraci\'on total del curso: 40 Hrs. (24 Hrs teor\'ia, 16 Hrs practica)
\end{flushleft}

\chapter {IMPORTANCIA DE LA PROTECCI\'ON RADIOL\'OICA.}
El uso de la radioactividad en la ingeniería es uno de los aspectos mas difundidos de la tecnología nuclear, puesto
que forma parte del pequeño numero de técnicas de control de calidad mediante pruebas no destructivas.\\

Pero así como son muchas las ventajas que nos ofrece también se corren muchos
riesgos si no sabemos emplearla bajo ciertos procedimientos que garanticen nuestra
protección contra la radiación.\\

Estar conscientes de que le programa de protección radiológica esta hecho para seguridad nuestra es dar el primer paso en la comprensión de este curso.\\

El objetivo principal de la protección radiológica es lograr que las operaciones controladas por la Dirección de un centro se desarrollen sin riesgo para la salud y la seguridad de las personas que estén dentro o fuera de el, así como la protección del medio ambiente, contra los riesgos que se derivan de las actividades que, por las características de ciertos materiales o equipos que se utilizan, pueden implicar la exposición a radiaciones ionizantes. Para conseguirlo, hay que aplicar los reglamentos que estén en vigor y todas las demás normas pertinentes, y vigilar debidamente las operaciones para comprobar la eficacia de las medidas de protección adoptadas.\\


Las responsabilidades que se adquieren al ser nombrado como personal ocupacionalmente expuesto son:
\begin{enumerate}
  \item Conocer y aplicar correctamente los principios b\'asicos de seguridad radiol\'ogica.
  \item Conocer el uso correcto de fuentes de radiaci\'on ionizante, equipo detector, y accesorios.
  \item Usar durante la jornada de trabajo el vestuario, equipo, accesorios y dispositivos, que es requieran de acuerdo al o establecido en le manual de seguridad radiol\'ogica y a lo dispuesto por el encargado de seguridad radiol\'ogica.
  \item Informar al encargado de seguridad radiol\'ogica sobre cualquier situación del alto riesgo, incidente o accidente radiol\'ogico.
\end{enumerate}

\chapter{ESTRUCTURA AT\'OMICA}

En la naturaleza, la materia se encuentra en una gran variedad de formas, Desde el punto de vista de la química, las sustancias puras se distinguen según sean compuestas o elementales. Los elementos son sustancias que no pueden descomponerse mas mediante métodos químicos.\\

La forma m\'inima en que se puede encontrar un compuesto se llama mol\'ecula y la forma m\'inima en que se puede encontrar un elemento se llama \ul{\'atomo}.\\

Aunque etimológicamente la palabra átomo significa sin corte, en el presente siglo se ha establecido que los átomos están constituidos por \ul{electrones} [$\mathrm{e^-}$] (partículas con carga eléctrica negativa que giran en orbitas especificas alrededor del núcleo.), y un núcleo el cual esta a su vez formado por \ul{protones} [$\mathrm{p^+}$] (partículas con carga eléctrica positiva) y \ul{neutrones} [$\mathrm{n^0}$] (partículas con carga eléctrica neutra).\\

Los átomos son generalmente neutros porque el numero de electrones que giran alrededor del núcleo con sus cargas negativas es igual al numero de protones en el núcleo con sus cargas positivas.\\

El numero de protones determina la identidad del elemento químico, cuando un mismo elemento tiene diferente numero de neutrones, en estos casos se habla de isótopos del elemento.\\

A cada elemento se le asocia con el numero de de orden que le corresponde dentro de la tabla periódica, a este numero se le conoce como \ul{Numero at\'omico} que corresponde al numero de \ul{protones} que posee el núcleo, y se representa con la \ul{letra A}.\\

Para diferenciar un isótopo de otros se usa la siguiente notación:
\newline

\begin{wrapfigure}{l}{.3\textwidth}
  \hspace{20mm}
  \begin{centering}
    \Huge${}^A_ZX_N$
  \end{centering}
  \caption{Simbologia isotopo}
  % \vspace{-6}
\end{wrapfigure}

\hspace{21mm}      donde: X= Simbolo quimico
\hspace{35mm}   A= Numero masico
\hspace{38mm}   = suma de protones y neutrones
\hspace{35mm}   Z= Numero Atomico
\hspace{35mm}   N= Numero de neutrones


% \hspace{-25}
Por ejemplo:

\begin{wrapfigure}{l}{.3\textwidth}

  \caption{Carbono-14}
  \begin{centering}
  \end{centering}
\end{wrapfigure}

\begin{wrapfigure}{l}{.3\textwidth}
  % \hspace{20mm}
  \begin{centering}
    \huge${}^A_ZX_N$
  \end{centering}
  \caption{Simbologia isotopo}
  % \vspace{-6}
\end{wrapfigure}

\hspace{21mm}      donde: X= Simbolo quimico
\hspace{35mm}   A= Numero masico

\hspace{38mm}   = suma de protones y neutrones

\hspace{35mm}   Z= Numero Atomico

\hspace{35mm}   N= Numero de neutrones

\chapter{RADIATIVIDAD Y RADIACION}

A medida que los nucleos se hacen mas complejos por el numero de particulas que
contienen el tamano de ellos se va incrementando y por ende las fuerzas
nucleares se van debilitando mientras qeu las electrostaticas se incrementan,
esto ocasiona que los nucleos pierdan estabilidad, y es por esta razon (aunque
no es la unica) es que la mayor parte de los nucleos pesados presentan
radiactividad.\\

La radiactividad se manifiesta a traves de la emision espontanea
de ciertas radiaciones por parte de los nucleos atomicos.Dichas
radiaciones son capaces de ¡impresionar placas fotograficas, de producir centelleos
o destellos luminosos al incidir en materiales fluorescentes, y de producir ionizacion en los materiales que atraviesan. Estas propiedades sirven de base para la operacion de los detectores de las mismas.\\

Una condicion para que la desintegracion de un nucieo en otros
dos ocurra espontaneamente es que la masa del nucleo original
sea suficiente para dar lugar a la formacion de los nucleos
producto y la energia de su movimiento, si esto no puede suceder
espontaneamente se dice que el nucleo es estable contra la desintegracion correspondiente.\\

Existen tres tipos de radiaciones emitidas por sustancias
radiactivas, a las cuales se les asignaron los nombres de
radiaciones alfa, beta y gamma en base a sus poderes de
penetracion crecientes, y a sus poderes de ionizacion
decrecientes, respectivamente.\\

Los rayos alfa: consisten de 2 protones y 2 neutrones unidos
formando un particula de gran masa con carga doblemente positiva
y con un gran poder de ionizacion, estos tienen muy corto alcance
y por lo tanto muy poco poder de penetracion.\\

Los rayos beta: Son similares a los electrones en masa y Carga,
tienen mayor poder de penetrcion que los alfa, estos requieren
de varias hojas de aluminio o varios metros de aire para ser
Catenidos por completo, y su poder de ionizacon es intermedio.\\

Los rayos gamma: son radiaciones electromagneticas de la misma naturaleza que la luz que siempre atraviesa cualquier espesor de material absorbedor, inclusive el plomo, y producen una ionizacion relativamente baja.\\

La actividad de una muestra radioactiva es el numero de desintegraciones que ocurren en la misma por unidad de tiempo, es decir que la actividad esla rapidez de desintegracion de la muestra.

Para medir la cantidad de material radioactivo se usa como unidad el Becquerel que se podria definir con la cantidad de material radioactivo que produce 1 desintegracion por segundo, sin emabargo la unidad más comunmente usada para medir actividad es el Curie, el cual se define com la actividad de una muestra en la que ocurren $3.7 x 10^{10}$ desintegracciones por segundo.\\

Algunos submúltiplos del curie son:

\begin{align*}
  1 \texbf{mCi} =& 1 \textbf{milicurie}  = 3.7 \times 10^{7} \frac{des}{s}\\
  1 \text{Ci}   =& 1 \textbf{microcurie} = 3.7 \times 10^{4} \frac{des}{s}\\
  1 \text{pCi}  =& 1 \textbf{picocurie}  = 3.7 \times 10^{-2} \frac{des}{s}
\end{align*}

La vida media $(T\frac{1}{2})$ o periodo de simidesintegracion corresponde al tiempo necesario para que la actividad se reduzca a la mitad de su valor inicial.\\

Para el calculo de la actividad de una fuente se pueden usar 3 metodos:

\section{Metodo Analitico}
Se realiza utilizando la siguiente formula:
\[A = Aoe^{\lambda t}\]

donde A es la actividad al tiempo (t), Ao es la actividad original al tiempo $t=0$, e es la base de los logaritmos neperianos = 2.71828, t es el tiempo transcurrido desde que se meidio o calculo la actividad original hasta el momento en que se desea calcular la nueva actividad y $\lambda$ es la constante de desintegracion = \frac{0.693}{\text{vida media}}\\

Supongamos que deseamos calcular la actividad que tendra el 20 de Marzo de 1995 una fuente de Ir-192 que se adquirio el 22 de Noviembre de 1994 con 110 Ci.
\begin{itemize}
  \item Primero se contabilizan los dias transcurridos: son 8 dias de noviembre + 31 de diciembre + 31 de enero + 29 de febrero + 20 de marzo = 118 dias.
  \item La vida media del Ir-192 es de 74 dias.
    \begin{align*}
      A=& Aoe^{\lambda t}\\
      A=& 110 e ^{-\frac{0.693 * 118}{74}}\\
      A=& 110 e^{1.105}\\
      A=& 110 * 0.33\\
      A=& 36.43 \text{Ci}
    \end{align*}
\end{itemize}

\section{Metodo tabular}
Consiste en encontrar el numero de dias transcurridos desde la fecha de calibracion de la fuente hasta la actual y buscar en la tabla (la cual es especifica para cada isotopo) el factor que corresponda. Luego lo multiplicamos por la actividad inicial y se obtiene el resultado. \\
Por ejemplo: Calcular la actividad actual de una fuente con una actividad inicial de 106 Ci, habiendo transcurrido 156 dias hasta la fecha de hoy.



\begin{center}
\begin{tabular}{ |c|c|c|c|c|c| }
 \hline
DIAS  &FACTOR    &DIAS  &FACTOR  &DIAS     &FACTOR\\
 \hline
25    &0.79      &200   &0.15    &375      &0.030\\
50    &0.63      &225   &0.12    &400      &0.024\\
75    &0.50      &250   &0.096   &425      &0,019\\
100   &0.39      &275   &0.076   &450      &0.015\\
125   &0.31      &300   &0.060   &475      &0.012\\
150   &0.25      &325   &0.048   &500      &0,009\\
175   &0.19      &350   &0.038   &525      &0.007\\
 \hline
\end{tabular}

TABLA DE DECAIMIENTO PARA Ir-192\\
\end{center}

Tomamos el factor que mas se aproxime a los 156 dias, en este caso es: 150 dias = 0.25\\
Actividad Actual = 0.25 \times 106 = 26.5 Ci

\Section{Metodo gráfico}
Consiste en graficar la actividad contra el tiempo tomando en cuenta que despues de cada vida media la actividad original va decayendo a la mitad. (Se utiliza papel similogaritmico para graficarlo), siendo la grafica exclusiva para el Ir-192.\\

\begin{center}
  \includegraphics[scale=0.5]{graficaIridio.png}
\end{center}

\chapter{INTERACCION DE LA RADIACION CON LA MATERIA}
Se entiende por radiacion ionizante, a aquellas ondas electromagneticas o particulas capaces de interactura con los electrones de los atomos que componen la materia, para producir iones.\\

Los iones son atomos o moleculas cargados electricamente.\\

Las ondas electromagneticas capaces de producir ionizacion, son: los rayos X y gamma. Las particulas que producen ionizacion son: las alfa, beta, protones y neutrones.\\

La ionizacion es el proceso por medio del cual un atomo o molecula adquiere una carga electrica positiva o negativa mediante la extraccion o adicion de electrones.\\

Excitacion es el estado donde se encuentra un excedente de energia que se elimina mediante la emision de rayos gamma.\\

Absorcion es la capacidad de un material de absorber o detener libres hidronio ($HO_{2}$) y peroxido de hidrogeno ($H_{2}O_{2}$), los cuales son biologicamente muy dañinos y pueden producir la muerte de la celula o puede mantenerla en estado latente donde se puede llegar a realizar una mitosis.\\

El nucleo es la parte mas sensible a la radiacion, a este nivel la radiacion gamma colisiona con los enlaces de la molecula de DNA, alterando el gene en el cual se encuentra la molecula de DNA, este gene alterado puede llegar a producir un efecto en los cromosomas, que al dividirse produce celulas alteradas con funciones diferentes a las de la celula madre; Esto se conoce como efecto genetico de la radiacion ionizante.\\

En general, podemos decir que la alteracion de la radiacion ionizante con las celulas, puede producir en ellas.
\begin{itemize}
  \item Aumento o disminucion de volumen.
  \item Su muerte
  \item Un estado latente
  \item Mutaciones en las celulas hijas
\end{itemize}

La Radiosensibilidad es la respuesta que las celulas ofrecen
a la accion de la radiacion, la cual es directamene proporcional
a la capacidad reproductiva de las celulas y varia inversamente
con el grado de diferenciacion de los mismos, asi como de el
estado metabolico, su estado de division y su estado de nutricion\\

Por lo, tanto las celulas que mas rapidamente se dividen son
las mas radiosensibles (hematopoyeticas, piel, cristalino,
epitelio, gastrointestinales, gonadas, etc.).\\

En funcion a los dos tipos de celulas que se encuentran en los
organis no, los efectos biologicos de la radiacion se clasifican
en dos tipos: somaticos y hereditarios.\\

Los somaticos son aquellos que se manifiestan en el individuo
que se ha expuesto a la radiacion y com su nombre lo indica,
se deben al daño recibido en sus celulas somaticas o diploides;
ademas este daño queda limitado solamente al individuo. Algunos
ejemplos de este tipo de alteraciones son la esterilidad,
disminucion de celulas en la medula osea e induccion al Cancer.\\

Los hereditarios o geneticos son el resultado del dano recibido
en las celulas gameticas y sus efectos se presentaran en la
descendencia de las celulas irradiadas. Un ejemplo serian los
procesos mutagenicos en el ADN, los cuales serian heredados
a sus descendientes.\\

Actualmente los efectos biologicos de la radiacion, de acuerdo
a su probabilidad de incidencia, se clasifican en dos tipos:
efectos deterministicos (no estocasicos) y efectos estocasticos.\\

Los efectos deterministicos se definen como aquellos que se
producen a partir de una dosis de umbral y aumentan en severidad con la dosis. El efecto se agudiza para dosis altas recibidas en un tiempo corto, estos efectos se derivan de la muerte celular.\\

La dosis umbral para algunos efectos deterministicos se presenta a continuacion, cuando la irradacion es a cuerpo total.\\

% TODO: Insertar tabla de efectos pagina 13

Los efectos estocasticos son aquellos que aparentemente no tiener
una dosis umbral a partir de la cual se manifiestan, le
probabilidad de que ocurra el efecto independientemente de su
severidad, se considera en funcion de la dosis, ademas debe
tenerse en cuenta que la mayoria de estos efectos exister
normalmente en poblaciones no irradiadas. Estos efectos se pueden
dividir en cuatro grupos:\\

\begin{itemize}
\item Induccion de leucemia.
\item Induccion de otros canceres.
\item Enfermedades hereditarias.
\item Acortamiento de la vida.
\end{itemize}

El proposito de la proteccion radiologica es evitar los efectos
deterministicos fijando limites de dosis por debajo de sus
umbrales.\\

Dentro de las formas de exposicion a la Radiacion, se encuentra
la irradiacion externa la que recibe el organismo desde fuentes
exteriores al mismo, los neutrones se consideran el tipo de
radiacion mas peligrosa, tanto por ser los mas dificiles de
detectar, como por su efecto sobre los tejidos, ya que al
interaccionar con los atoms de hidrogeno producen protones cuyo
elevado poder de ionizacion causa un daño muy localizado.\\

La irradiacion interna se presenta cuando la fuente se encuentra
dento del organismo. Los materiales radiactivos pueden ingresar
al cuerpo por ingestion, por inhalacion, por absorcion a traves
de la piel o bien a la sangre por algun corte o herida. El dano
dependera del tipo de radiacion emitida, su energia, la vida
media del radionuclido, su permanencia en le organismo y el
organo en que se localice. En este Caso los emisores alfa y
beta son los mas peligroso.\\

A continuacion se anexan cuadros que resumen los efectos de
la irradicion al organismo.\\

% TODO: Insertar tabla de efectos pagina 15

Como ya se menciono la utilizacion de materiales radiactivos
implica la presencia de alguno o varios de los riesgos
siguientes:
%TODO: formatear correctamente
Irradiacion
Externa (gamma, beta)
Interna (alfa, beta, gamma)
Riesgos:
Contaminacion Externa (gamma, beta)
Interna (alfa, beta, gamma)

Los materiales radiactivos podran liberarse al medio ambiente
cuando ocurra algun accidente que afecte la integidad del
contenedor y capsula. Cuando esto ocurre, la presencia del
material radiactivo ocasiona riesgos de contaminacion externa
y/o interna ya que el material puede embarrarse en la piel
durante las actividades de trabajo e incluso ingerirse o
inhalarse, dependiendo de las caracteristicas fisicas del
material radiactivo. \\

A continuacion se mencionan algunas medidas de proteccion contra
la irradicion externa que se deberan seguir siempre al tomar
una radiografia
\begin{itemize}
\item Antes de empezar el trabajo checar que su dosimetro de lectura
directa no se encuentra saturado o fuera de escala y que
su dosimetro de pelicula esta en buen estado.

\item Verificar que el detector giger tiene sus baterias en buen
estado y que su alarma este encendida.

\item Antes de empezar la irradiacion se debe acordonar una area
aproximada de 15 mts. de la fuente con letreros que indiquen
la pesenicia de radiacion.

\item Durante el tiempo que dure la exposicion se debe mantener
alejado de la fuente pero controlando el area para evitar
que algunas personas pudieran acercarse y recibir una
exposicion innecesariamente.

\item Al terminar la exposicion se debe verificar con el detector
geiger que la fuente efectivamente regreso al contenedor
y que se encuentra segura. Esto se hara tomando los niveles
de radiaicion cerca del contendor y de la manguera
posicionadora.

\item Volvera a leer su dosimetro de bolsillo y tomara nota de
la dosis que recibio durante su jornada de ese dia. Al final
del dia registra su dosis en las formas ya.  previstas para
ello y al final de cada semana las entregara al encargado.

\item En caso de exposiciones con grado de dificultad mayor, como
por ejemplo en tuberias aereas, en donde el tecnico debe
permanecer cerca de la fuente ya que el tiempo de irradiacion
va a ser muy corto y no le permite alejarse y regresar, el
uso de el colimador es inidispensable ademas de que la persona
tendra que planear el trabajo en forma tal que se apliqur
lo que hasta aqui ha aprendido.
\end{itemize}




\end{document}
